\chapter{Wie funktioniert Ausrüstung?}

\section{Was ist Ausrüstung?}

Ausrüstung sind Gegenstände, die ein Held bei sich trägt und verwenden kann.

Ausrüstung wird durch Karten dargestellt. Eine Ausrüstungskarte kann zum Beispiel ein Schwert, einen Schild, eine Rüstung, einen Helm oder einen besonderen Gegenstand zeigen.

Auf jeder Ausrüstungskarte steht, wo der Gegenstand ausgerüstet wird und welche Wirkung er im Spiel hat.

\section{Ausrüstungsplätze}

Viele Ausrüstungskarten haben oben links ein Symbol. Dieses Symbol zeigt, welchen Ausrüstungsplatz der Gegenstand benötigt.

Ein Held hat folgende Ausrüstungsplätze:

\begin{itemize}
  \item Kopf
  \item zwei Hände
  \item Körper
  \item Gürtel
  \item Füße
\end{itemize}

Ein Gegenstand gilt nur dann als ausgerüstet, wenn er an seinem passenden Ausrüstungsplatz liegt. Nur ausgerüstete Gegenstände können verwendet werden.

An Kopf, Körper, Gürtel und Füßen kann jeweils nur ein Gegenstand ausgerüstet werden. Für die Hände gelten besondere Regeln.

\subsection{Hände}

Die Hände funktionieren etwas anders als die übrigen Ausrüstungsplätze.

Ein Held hat zwei Hände. Manche Gegenstände benötigen eine Hand, andere benötigen beide Hände. Ein Held kann Gegenstände für die Hände ausrüsten, die zusammen höchstens zwei Hände benötigen.

Ein Schwert benötigt zum Beispiel eine Hand. Ein Schild benötigt ebenfalls eine Hand. Ein Zweihänder benötigt beide Hände.

\subsection{Verstaute Ausrüstung}

Ein Held kann auch weitere Ausrüstung mit sich tragen.

Diese Ausrüstung ist verstaut, zum Beispiel im Rucksack, in einer Tasche oder am Gepäck. Verstaute Ausrüstung gilt nicht als ausgerüstet und kann deshalb nicht verwendet werden.

Nur Ausrüstung, die an ihrem passenden Ausrüstungsplatz liegt, gilt als ausgerüstet.

\subsection{Kleine Gegenstände}

Neben seiner Ausrüstung kann ein Held beliebig viele kleine Gegenstände mit sich tragen.

Kleine Gegenstände benötigen keinen Ausrüstungsplatz. Sie können im Rucksack liegen, in Taschen stecken oder als Schmuck am Körper getragen werden.

\section{Wie liest man eine Ausrüstungskarte?}

Je nach Art der Ausrüstung stehen unterschiedliche Informationen auf einer Karte.

Nicht jede Ausrüstungskarte enthält alle der folgenden Angaben.

\subsection{Angaben zu den Angriffsoptionen}

Waffen enthalten eine oder mehrere Angriffsoptionen. Eine Angriffsoption beschreibt eine bestimmte Art, die Waffe einzusetzen.

Jede Angriffsoption legt fest:

\begin{itemize}
  \item die Angriffsfertigkeit, mit der angegriffen wird. Die Angriffsfertigkeit bestimmt auch, welche Würfel für die Angriffsprobe verwendet werden und welche Gegner in Reichweite sind,
  \item die Schadensformel, mit der Sterne in Schadenswürfel umgerechnet werden,
  \item mögliche besondere Effekte.
\end{itemize}

Jede Angriffsoption steht in einer eigenen Zeile.

\subsection{Angaben zum Ausweichen}

Manche Ausrüstung, meist Rüstungen, ermöglicht dem Träger eine Ausweichprobe. Mit dieser kann ein Held dem Angriff eines Gegners teilweise oder sogar vollständig ausweichen.

Leichte Rüstungen erlauben meist bessere Ausweichproben, während sehr schwere Rüstungen gar kein Ausweichen mehr erlauben.

Trägt ein Held gar keine Rüstung, sein Körper-Ausrüstungsplatz ist also leer, darf er als Ausweichprobe eine Turnen-Probe durchführen.

Erlaubt die Ausrüstung eine Ausweichprobe, gibt sie an, welche Attribute dafür verwendet werden können:

\begin{itemize}
  \item Eine leichte, agile Rüstung verwendet eventuell INT und GEW: Die Rüstung ermöglicht es dem Träger, geschickt zur Seite zu springen.
  \item Eine magische Ausrüstung verwendet vielleicht CHA und INT: Ein magischer Zauber leitet den Schlag eines Gegners um.
  \item Ein Parierstab verwendet vielleicht KRA und KON: Der Held hält aktiv den Schlag auf.
\end{itemize}

Wie bei allen Proben wird bei einer Ausweichprobe auch der Heldenwürfel gewürfelt.

Besonders gute Rüstungen stellen dem Helden eventuell zusätzliche Würfel für Ausweichproben zur Verfügung.

Steht auf einer Ausrüstung \emph{Ausweichen: Heldenwürfel + Attribut 1 + Attribut 2}, bedeutet dies, dass diese Ausrüstung eine Ausweichprobe erlaubt.

Steht auf einer Ausrüstung \emph{Ausweichen: +x Würfel}, bedeutet dies, dass ein Held beim Ausweichen zusätzliche Würfel erhält. Diese Würfel sind für gewöhnlich blau und werden unabhängig von Attributs- oder Fertigkeitswerten zum Würfelpool hinzugefügt.

\paragraph{Beispiel}
Gadhra trägt eine Kettenrüstung und einen kleinen Deckel:

\begin{itemize}
  \item Auf der Kettenrüstung steht \emph{Ausweichen: Heldenwürfel + INT + GEW}.
  \item Auf dem kleinen Deckel steht \emph{Ausweichen: +1 blauer Würfel}.
\end{itemize}

Das bedeutet: Wenn der Held ausweicht, würfelt er seinen Heldenwürfel, seine Würfel für Intuition und Gewandtheit und einen zusätzlichen blauen Würfel.

\subsection{Besondere Effekte}

Einige Ausrüstungsgegenstände besitzen zusätzlich besondere Effekte.

Diese stehen als Regeltext oder als Symbole auf der Karte und gelten genau so, wie sie dort beschrieben sind.

Einige Effekte gelten dauerhaft.

Andere können nur unter bestimmten Bedingungen genutzt werden oder müssen bezahlt werden, zum Beispiel durch Mondkristalle.

Manche Effekte fordern den Spieler auf, die Ausrüstungskarte umzudrehen. Die Karte gilt dann als erschöpft. Ein Effekt kann erst wieder genutzt werden, wenn die Karte wieder bereitgemacht wurde. Solange eine Karte erschöpft ist, gilt sie weiterhin als ausgerüstet und belegt ihren Ausrüstungsplatz. Ihre Angaben und Effekte können jedoch nicht genutzt werden.

Erschöpfte Karten werden spätestens bei der nächsten Rast wieder aufgedeckt. Durch besondere Effekte, Zaubersprüche oder erzählerisch passende Handlungen können Karten auch früher bereitgemacht werden, wenn die Spielleitung dies erlaubt.

\section{Weiterführende Informationen}

Im Anhang findet sich eine Liste aller Ausrüstungsgegenstände sowie eine Übersicht über alle Symbole, die auf Ausrüstungskarten vorkommen können, einschließlich ihrer Erklärung.
